\documentclass[11pt]{article}
\usepackage{fontspec}
\usepackage[french,english]{babel}
\usepackage{amsmath}
\usepackage{amssymb}
\usepackage{amsthm}
\usepackage{mathrsfs}
\usepackage{tikz-cd}
\usepackage{enumitem}
\usepackage{lmodern}

\newcommand{\printedpage}[1]{\marginpar{\footnotesize p.~#1}}
\newcommand{\msp}[1]{\marginpar{\footnotesize [#1]}}
\newcommand{\cX}{\mathscr{X}}
\newcommand{\Ob}{\operatorname{Ob}}

\begin{document}

\printedpage{150}
% head: A. Grothendieck : Dérivateurs, Ch. XII, Édition des Universités de Montpellier II et Paris VII
% foot: Version 30 mars 2003

It remains to prove c). I am going to deduce it from a

\textbf{Proposition 17.} \textit{Let there be a commutative square}
\[
Q \quad = \quad
\begin{tikzcd}[ampersand replacement=\&]
X_0 \arrow[r, "i_1"] \arrow[d, "i_2"'] \& X_1 \arrow[d, "f_1"] \\
X_2 \arrow[r, "f_2"] \& Y
\end{tikzcd}
\]
\textit{in} Cat\textit{, and let $\cX_Q$ be the fibered category over}
\[
\tilde{\Phi} \quad (\simeq \mathbf{\Delta}^1 \times \mathbf{\Delta}^1) \quad = \quad
\begin{tikzcd}[ampersand replacement=\&]
0 \& 1 \arrow[l] \\
2 \arrow[u] \& 3 \arrow[l] \arrow[u]
\end{tikzcd}
\]
\textit{which integrates it. Let $\cX$ be the inverse image of the closed subset}
\[
\begin{tikzcd}[ampersand replacement=\&]
0 \& 1 \arrow[l] \\
2 \arrow[u] \&
\end{tikzcd}
\]
\msp{page 235}
\textit{in $\cX_Q$, which integrates the diagram}
\[
\begin{tikzcd}[ampersand replacement=\&, row sep=small]
 \& X_1 \\
X_0 \arrow[ru] \arrow[rd] \& \\
 \& X_2
\end{tikzcd}
\]
\textit{. For $Q$ to be $W$-cocartesian, it is necessary and sufficient that the inclusion}
\[
i : \cX \longrightarrow \cX_Q
\]
\textit{be in $W$.}

Since $\cX_Q$ is smooth over $\tilde{\Phi}$, and 3 is an initial object of $\tilde{\Phi}$, it follows that the inclusion of the fiber of $\cX_Q$ at the vertex 3 of $\tilde{\Phi}$ is in $W$,
\[
Y \xrightarrow{i_Y} \cX_Q \qquad\qquad \in W.
\]

Consider then the diagram
\[
\begin{tikzcd}[ampersand replacement=\&, column sep=small]
\cX \arrow[rr, "\varphi_Q"] \arrow[rd, "i"', ""{name=direct,pos=.55,above right}] \& \& Y \arrow[ld, "i_Y\in W", ""{name=via,pos=.30,below left}] \\
 \& \cX_Q. \& \arrow[from=via,to=direct,dotted,shorten <=1pt,shorten >=1pt,"\alpha"']
\end{tikzcd}
\]

\printedpage{151}
% head: A. Grothendieck : Dérivateurs, Ch. XII, Édition des Universités de Montpellier II et Paris VII
% foot: Version 30 mars 2003

This diagram is not commutative, but there is a commutativity arrow
\[
\alpha : i_Y\varphi_Q \longrightarrow i.
\]

We know that this implies that
\[
i \in W \qquad\qquad \Longleftrightarrow \qquad\qquad (i_Y\varphi_Q) \in W
\]
(cf. lemma p. 2), and since $i_Y \in W$,
\[
(i_Y\varphi_Q) \in W \qquad\qquad \Longleftrightarrow \qquad\qquad \varphi_Q \in W,
\]
whence
\[
i \in W \qquad\qquad \Longleftrightarrow \qquad\qquad \varphi_Q \in W,
\]

\hfill q.e.d.

We can now prove

\msp{page 236}

\textbf{Proposition 18} (\footnote[127]{should be true in every category of models having enough cofibrations.}). (Assuming W(7 bis).) \textit{Consider a commutative diagram}
\[
\begin{tikzcd}[ampersand replacement=\&]
X_0 \arrow[r] \arrow[d] \arrow[rd, phantom, "1"] \& X_1 \arrow[r] \arrow[d] \arrow[rd, phantom, "2"] \& X_2 \arrow[d] \\
X'_0 \arrow[r] \& X'_1 \arrow[r] \& X'_2
\end{tikzcd}
\tag*{$(\ast)$}
\]
\textit{whence a composite square}
\[
\begin{tikzcd}[ampersand replacement=\&]
X_0 \arrow[r] \arrow[d] \arrow[rd, phantom, "3"] \& X_2 \arrow[d] \\
X'_0 \arrow[r] \& X'_2.
\end{tikzcd}
\]
\textit{Suppose that square} 1 \textit{is $W$-cocartesian. Then square} 2 \textit{is $W$-cocartesian if and only if square} 3 \textit{is.}

\textsc{Proof.}

\texttt{[Now, }$X_i$\texttt{ and }$X'_i$\texttt{ exchange their roles.]}

\printedpage{152}
% head: A. Grothendieck : Dérivateurs, Ch. XII, Édition des Universités de Montpellier II et Paris VII
% foot: Version 30 mars 2003

Let
\[
I \quad = \quad
\left(
\begin{tikzcd}[ampersand replacement=\&]
0 \arrow[d] \& 1 \arrow[l] \arrow[d] \& 2 \arrow[l] \arrow[d] \\
0' \& 1' \arrow[l] \& 2' \arrow[l]
\end{tikzcd}
\right)
\]
and let $\cX$ be the integral of the diagram $(\ast)$, a fibered category over $\cX$. We denote by $\cX_{010'}$ the full subcategory of $\cX$ which is the inverse image of the full subcategory % [sic] fibrée sur $\cX$ -> fibrée sur $I$, a category fibred over itself; the base is the category $I$ just defined
\[
\underbrace{I_{010'}}_{=\{0,1,0'\}} \ \subseteq I,
\]
and likewise for every other subset $A$ of $\Ob I = \{0,1,2,0',1',2'\}$, giving a full subcategory $\cX_A$ of $\cX$. We put
\[
U \quad = \quad
\begin{tikzcd}[ampersand replacement=\&]
0 \arrow[d] \& 1 \arrow[l] \& 2 \arrow[l] \\
0' \& \&
\end{tikzcd}
, \qquad
V \quad = \quad
\begin{tikzcd}[ampersand replacement=\&]
0 \arrow[d] \& 1 \arrow[l] \arrow[d] \& 2 \arrow[l] \\
0' \& 1' \arrow[l] \&
\end{tikzcd}
,
\]
whence
\[
\cX_U = \cX_{0120'}, \qquad\qquad \cX_V = \cX_{0120'1'}.
\]
The arrows between the $\cX_A$ are the inclusions

\msp{page 237}

deduced from inclusions $A \hookrightarrow A'$ in $\mathfrak{P}(\Ob I)$. Thus, the $W$-cocartesian character of squares 1, 2, 3 is expressed respectively by the condition that the inclusions
\[
\begin{aligned}
\cX_{010'} \quad &\xrightarrow{\varphi_1} \quad \cX_{010'1'} \\
\cX_{121'} \quad &\xrightarrow{\varphi_2} \quad \cX_{121'2'} \\
\cX_{020'} \quad &\xrightarrow{\varphi_1} \quad \cX_{020'2'} % [sic] $\varphi_1$ -> $\varphi_3$, this inclusion is called $\varphi_3$ in the diagram that follows
\end{aligned}
\]
be in $W$. Consider then the diagram of inclusion arrows
\printedpage{153}
% head: A. Grothendieck : Dérivateurs, Ch. XII, Édition des Universités de Montpellier II et Paris VII
% foot: Version 30 mars 2003
\[
\begin{tikzcd}[ampersand replacement=\&]
\cX_{020'} \arrow[r, "\varphi_3"] \arrow[d, "\alpha"'] \& \cX_{020'2'} \arrow[d, "\beta"] \\
\cX_U \arrow[r, "\psi_3"] \arrow[d, "\psi_1"'] \& \cX \arrow[d, "\mathrm{id}"] \\
\cX_V \arrow[r, "\psi_2"] \& \cX \\
\cX_{121'} \arrow[r, "\varphi_2"] \arrow[u, "\alpha'"] \& \cX_{121'2'}. \arrow[u, "\beta'"']
\end{tikzcd}
\]

I say that in this diagram we have
\[
\alpha, \beta, \alpha', \beta' \in W
\tag*{$(\ast)$}
\]
and
\[
\varphi_1 \in W \ \ \text{(i.e. square 1 $W$-cocartesian)} \qquad\qquad \Longleftrightarrow \qquad\qquad \psi_1 \in W.
\tag*{$(\ast\ast)$}
\]

The relations $(\ast)$ already imply that % [sic] déja -> déjà, missing grave accent
\[
\begin{aligned}
\varphi_2 \in W \quad &\Longleftrightarrow \quad \psi_2 \in W \\
\varphi_3 \in W \quad &\Longleftrightarrow \quad \psi_3 \in W,
\end{aligned}
\]
and under the $W$-cocartesian hypothesis $(\ast\ast)$, implying $\psi_1 \in W$,

\msp{page 238}

we find
\[
\psi_2 \in W \qquad\qquad \Longleftrightarrow \qquad\qquad \psi_3 \in W,
\]
whence
\[
\varphi_2 \in W \qquad\qquad \Longleftrightarrow \qquad\qquad \varphi_3 \in W,
\]
which is what was to be proved.

It remains to establish $(\ast)$ and $(\ast\ast)$. The relations $(\ast)$ come from the observation that if $i : A \longrightarrow A'$ is an inclusion of full subcategories of $I$, such that $i$ is $W$-\textit{aspheric} (i.e. $A/j$ $W$-aspheric for $j \in \Ob A'\backslash\Ob A$), then $\cX_A \longrightarrow \cX_{A'}$ is so too (smooth base change), hence is \texttt{[in]} $W$. We must verify $W$-asphericity for the four inclusions
\[
\begin{array}{rcll}
\{0,2,0'\} & \longrightarrow & U = \{0,1,2,0'\} & \qquad\qquad \mathbf{\Delta}^0 \\[1ex]
\{0,2,0',2'\} & \longrightarrow & I & \qquad\qquad \mathbf{\Delta}^0/\mathbf{\Delta}^1 \\[1ex]
\{1,2,1'\} & \longrightarrow & V = \{0,1,2,0',1'\} & \qquad\qquad \mathbf{\Delta}^1,\ \Phi \\[1ex]
\{1,2,1',2'\} & \longrightarrow & I & \qquad\qquad \mathbf{\Delta}^1,\ \mathbf{\Delta}^1 \times \mathbf{\Delta}^1.
\end{array}
\]

\printedpage{154}
% head: A. Grothendieck : Dérivateurs, Ch. XII, Édition des Universités de Montpellier II et Paris VII
% foot: Version 30 mars 2003

(NB I have marked in the margin the categories we obtain as $A/j$, up to isomorphism, in the four cases considered - they all have an initial object, hence are $W$-aspheric.)

It remains to prove $(\ast\ast)$. Let
\[
U' = \{0,1,0',1'\}, \qquad\qquad \text{hence } V = U \cup U',\ \ U \cap U' = \{0,1,0'\},
\]
where $U$, $U'$ are \textit{closed subsets} of $V$ which cover $V$. Thus we will likewise have
\[
\cX_V = \cX_U \cup \cX_{U'}, \qquad\qquad \cX_U \cap \cX_{U'} = \cX_{010'},
\]

\msp{page 239}

\[
\begin{tikzcd}[ampersand replacement=\&]
\cX_{U\cap U'} \arrow[r] \arrow[d, "\varphi_1"] \& \cX_U \arrow[d, "\psi_1"] \\
\cX_{U'} \arrow[r] \& \cX_V.
\end{tikzcd}
\]

Thus by W(7 bis) (dual case of closed immersions), we have
\[
\varphi_1 \in W \qquad\qquad \Longleftrightarrow \qquad\qquad \psi_1 \in W,
\]

\hfill q.e.d.

Let us note for the record the corollary which had motivated us.

\textbf{Corollary.} \textit{Assume} W(7 bis) (as throughout this section). \textit{Let there be a cocartesian square in} Cat
\[
\begin{tikzcd}[ampersand replacement=\&]
X_0 \arrow[r, "i_1"] \arrow[d, "i_2"'] \& X_1 \arrow[d] \\
X_2 \arrow[r] \& X_3,
\end{tikzcd}
\]
\textit{with $i_1$ (or $i_2$) a $W$-cofibration (e.g. a Dwyer immersion). Then this square is $W$-cocartesian.}

This digression (p. 232 ff.) is not necessary, for that matter, to formulate right now the conclusion of our cogitations on $W$-homotopic cofibers: \textit{There is such a theory in} Cat, provided only that $W$ satisfies W(7 bis). Indeed, everything follows from the fact that every arrow $f : X \longrightarrow Y$ in Cat factors as
\[
X \xrightarrow{i_f} \tilde{Y} \xrightarrow{p_f} Y,
\]
with $i_f$ a $W$-cofibration (and even an imm-

\end{document}
